\documentclass[11pt]{article}
\usepackage[margin=1in]{geometry}
\usepackage{amsmath,amssymb}
\title{Report on ``Matrix Theory Reloaded: A BPS Road to Holography''\\
\small Blair, Lahnsteiner, Obers, Yan}
\date{}
\begin{document}
\maketitle

\section*{1. Principal achievement}
The paper offers a single organising principle for three subjects that are
usually treated separately: matrix theories, holography, and non-Lorentzian
string/M-theory. The key observation is that the decoupling limits giving
matrix theories on D$p$-branes (the ``M$p$T'' limits) are \emph{BPS} limits.
Applied to the curved brane background, the same limit automatically produces
the Maldacena near-horizon limit. When uplifted to M-theory, each such limit is
a DLCQ. This leads to the conjecture that holography is a
DLCQ$^n$/DLCQ$^m$ correspondence ($m>n$). In this language AdS$_5$/CFT$_4$ is
DLCQ$^0$/DLCQ$^1$, D1--D5 gives an example with $m>n+1$, and applying two
limits in sequence gives new dualities with non-Lorentzian bulks. A second,
more concrete result comes from reading off how the brane actions depend on the
decoupling parameter $\lambda=\omega^{-2}$. This gives explicit $T\bar T$-like
flow equations for the D$p$-brane DBI action and the M2-brane. The
three-dimensional M2 flow,
$\partial_\lambda\mathcal{L}=\tfrac14\sqrt{-\det\tau}\,[\mathrm{tr}\,\mathcal{T}^2-(\mathrm{tr}\,\mathcal{T})^2+2\lambda\det\mathcal{T}]$,
is especially simple. In the bulk, undoing the near-horizon limit then appears
as a deformation that generates the brane geometry out of the flat
non-Lorentzian geometry at asymptotic infinity. This is a fresh view of why the
bulk exists.

\section*{2. Principal drawback}
The central conjecture is mostly a reorganisation of known dualities and is not
yet a predictive statement. The paper does not give a dynamical test that would
separate DLCQ$^n$/DLCQ$^m$ from the standard near-horizon argument. The new
non-Lorentzian holographic examples are the least secure part. Their bulk
geometries come from an ad hoc extra rescaling of the harmonic function, which
amounts to rescaling $N$, and its physical meaning is unclear. The authors
also say openly that they have not checked that these geometries solve the
equations of motion and torsional constraints of M$p$T supergravity. That
theory has not itself been formulated, and the symmetries, charges and
supersymmetry of the solutions have not been worked out. The geometric picture
at asymptotic infinity has further weak points. Quantum corrections may
dominate there, the dilaton runs to strong coupling for $p>3$, and the
non-Lorentzian structure may depend on the choice of coordinates. Finally, the
$T\bar T$ flows are derived only for bosonic, abelian actions, while the claimed
link to matrix theory needs the non-abelian, supersymmetric case.

\section*{3. Most worthwhile follow-up}
Balancing feasibility and impact, I recommend \emph{extending the $p$-brane
$T\bar T$ flow equations to the non-abelian (and ideally supersymmetric) case
and testing them on matrix theory.} A concrete programme:
\begin{itemize}
\item Apply the flow, e.g.\ the M2 equation or its D$p$ analogue, to
U$(N)$ SYM as the initial condition. Solve it order by order in $\lambda$.
Compare the result with the known $\alpha'$ expansion of the non-abelian DBI
action: the symmetrised trace at $F^4$, and the known commutator ambiguities at
$F^6$ and beyond.
\item Use the M$p$T parametrisation of Myers' action as a target, and check
whether the flow fixes the ordering ambiguities.
\item For $p=2$, apply the flow to BLG/ABJM to look for an M-theoretic
non-abelian ``DBI'' endpoint.
\end{itemize}
This is feasible because it is a well-posed classical field-theory computation
that builds directly on equations the paper has already derived. The
$F^4$ term gives an immediate check. The potential impact is large. A successful
flow would turn the paper's main conjecture into something testable: undoing
the DLCQ/near-horizon limit would become a solvable deformation of matrix
theory, and it might fix the long-standing non-abelian DBI ordering problem.
If the flow fails to reproduce the $F^4$ term, that would be equally
informative about how far the framework can be trusted. A more ambitious but
riskier alternative is to formulate M$p$T supergravity and verify the
non-Lorentzian bulk solutions.

\end{document}
